% ============================================================================
% LANGENTWURF LE-02c: Possessivbegleiter
% Spanisch Klasse 7 - Sequenz 2: Mi familia
% ============================================================================
% Tier: 1 (vollständig)
% Doppelstunde: 2c (Std. 5-6 von 12)
% Fachfarbe: #c41e3a (Karminrot)
% ============================================================================

\documentclass[11pt,a4paper]{article}

\usepackage[utf8]{inputenc}
\usepackage[T1]{fontenc}
\usepackage[ngerman]{babel}
\usepackage{geometry}
\usepackage{fancyhdr}
\usepackage{lastpage}
\usepackage{xcolor}
\usepackage{tcolorbox}
\usepackage{enumitem}
\usepackage{tabularx}
\usepackage{longtable}
\usepackage{array}
\usepackage{booktabs}
\usepackage{amssymb}

% ============================================================================
% KONFIGURATION
% ============================================================================

\newcommand{\fach}{Spanisch}
\newcommand{\fachfarbe}{c41e3a}
\newcommand{\jahrgangsstufe}{7}
\newcommand{\schulart}{Gesamtschule}
\newcommand{\stundenthema}{Possessivbegleiter}
\newcommand{\reihenthema}{Mi familia}
\newcommand{\stundennummer}{2c}
\newcommand{\reihenstunden}{6}
\newcommand{\gession}{A1}
\newcommand{\lehrwerk}{Qué pasa 1}
\newcommand{\lehrwerkseite}{S. 64-66}
\newcommand{\lerngruppengroesse}{24}

% ============================================================================
% LAYOUT
% ============================================================================
\geometry{left=2.5cm, right=2.5cm, top=2.5cm, bottom=2.5cm}
\definecolor{fach}{HTML}{\fachfarbe}
\definecolor{gray}{HTML}{666666}
\definecolor{lightgray}{HTML}{f5f5f5}
\definecolor{successgreen}{HTML}{2a7a4b}
\definecolor{warningorange}{HTML}{b35c00}
\definecolor{basis}{HTML}{2a7a4b}
\definecolor{standard}{HTML}{2c5aa0}
\definecolor{erweiterung}{HTML}{7b1fa2}

\tcbuselibrary{skins,breakable}

\newtcolorbox{infobox}[1][]{
    colback=lightgray,
    colframe=fach,
    fonttitle=\bfseries,
    title=#1,
    breakable
}

\newtcolorbox{scorebox}{
    colback=white,
    colframe=fach,
    fonttitle=\bfseries,
    title=Didaktik-Score,
    breakable
}

\pagestyle{fancy}
\fancyhf{}
\fancyhead[L]{\small\textcolor{gray}{\fach{} -- Jg. \jahrgangsstufe{} -- \stundenthema}}
\fancyhead[R]{\small\textcolor{gray}{LE-\stundennummer{} | Tier 1}}
\fancyfoot[C]{\small\thepage{} / \pageref{LastPage}}
\renewcommand{\headrulewidth}{0.4pt}

% Differenzierungsmarker
\newcommand{\B}{\textcolor{basis}{\textbf{[B]}}}
\newcommand{\Std}{\textcolor{standard}{\textbf{[S]}}}
\newcommand{\E}{\textcolor{erweiterung}{\textbf{[E]}}}

% ============================================================================
% DOKUMENT
% ============================================================================
\begin{document}

% ============================================================================
% DECKBLATT
% ============================================================================
\begin{titlepage}
    \centering
    \vspace*{2cm}
    {\large\textcolor{gray}{\schulart}}\\[2cm]
    {\Huge\textcolor{fach}{\textbf{Langentwurf}}}\\[0.5cm]
    {\large LE-\stundennummer{} | Tier 1 (vollständig)}\\[2cm]
    {\LARGE\textbf{\stundenthema}}\\[0.5cm]
    {\large Sequenz: \reihenthema}\\[2cm]
    \begin{tabular}{rl}
        \textbf{Fach:} & \fach{} (A1) \\[0.3cm]
        \textbf{Klasse:} & \jahrgangsstufe \\[0.3cm]
        \textbf{Lehrwerk:} & \lehrwerk, \lehrwerkseite \\[0.3cm]
        \textbf{Doppelstunde:} & \stundennummer{} (Std. 5-6 von 12) \\[0.3cm]
    \end{tabular}
    \vfill
\end{titlepage}

\tableofcontents
\newpage

% ============================================================================
% 1. BEDINGUNGSANALYSE
% ============================================================================
\section{Bedingungsanalyse}

\subsection{Lerngruppe}
\begin{tabular}{ll}
    Kursgröße: & \lerngruppengroesse{} SuS \\
    Jahrgangsstufe: & \jahrgangsstufe \\
    Sprachniveau: & A1 (GER) -- Anfänger \\
    Fremdsprache: & 2. Fremdsprache (nach Englisch) \\
\end{tabular}

\subsection{Vorwissen aus 02a/02b}
\begin{itemize}
    \item Familienvokabular: padre, madre, hermano/a, abuelo/a
    \item Verb \textit{tener} (tengo, tienes, tiene)
    \item Verb \textit{ser} zur Personenbeschreibung
    \item Fragen: ¿Cómo se llama...? ¿Cuántos años tiene...?
    \item Bestimmte Artikel (el, la, los, las)
\end{itemize}

\subsection{Differenzierung (intern)}
\begin{tabular}{|l|p{3.5cm}|p{3.5cm}|p{3.5cm}|}
\hline
& \B{} \textbf{Basis} & \Std{} \textbf{Standard} & \E{} \textbf{Erweiterung} \\
\hline
\textbf{Ziel} & BR: Rezeption + Einsetzen & MR: Produktion mit Hilfe & GY: Freie Produktion \\
\hline
\textbf{Possessiv} & mi/tu/su erkennen & mi/mis, tu/tus, su/sus anwenden & Alle Formen + nuestro \\
\hline
\textbf{Anwendung} & Lückentext ausfüllen & Sätze bilden & Freie Familien\-beschreibung \\
\hline
\end{tabular}

% ============================================================================
% 2. SACHANALYSE
% ============================================================================
\section{Sachanalyse}

\subsection{Sprachliche Strukturen}

\textbf{Possessivbegleiter (betont, vorangestellt):}
\begin{center}
\begin{tabular}{|l|c|c|}
\hline
\textbf{Person} & \textbf{Singular} & \textbf{Plural} \\
\hline
yo & mi & mis \\
\hline
tú & tu & tus \\
\hline
él/ella/usted & su & sus \\
\hline
\end{tabular}
\end{center}

\textbf{Regeln:}
\begin{itemize}
    \item Der Possessivbegleiter richtet sich nach dem \textbf{Bezugswort} (nicht nach dem Besitzer!)
    \item Singular: mi madre, tu padre, su hermano
    \item Plural: mis hermanos, tus abuelos, sus primos
    \item Kein Unterschied maskulin/feminin bei mi/tu/su
    \item ACHTUNG: ``su'' kann bedeuten: sein/ihr/Ihr (usted)
\end{itemize}

\textbf{Redemittel:}
\begin{itemize}
    \item Mi madre se llama... (Meine Mutter heißt...)
    \item ¿Cómo se llama tu hermano? (Wie heißt dein Bruder?)
    \item Sus abuelos viven en... (Seine/Ihre Großeltern leben in...)
\end{itemize}

\subsection{Kontrastive Betrachtung}
\begin{itemize}
    \item \textbf{Positiv:} Possessiv-Konzept aus Deutsch/Englisch bekannt
    \item \textbf{Interferenz:} Im Deutschen: mein/meine (Genus); Spanisch: mi (kein Genus)
    \item \textbf{Interferenz:} Plural-Bildung: mis $\neq$ *meine (Numerus des Bezugsworts!)
    \item \textbf{Englisch-Transfer:} my/your/his $\rightarrow$ mi/tu/su (ähnliche Struktur)
\end{itemize}

% ============================================================================
% 3. DIDAKTISCHE ANALYSE
% ============================================================================
\section{Didaktische Analyse}

\subsection{Stellung in der Sequenz}
Dies ist die \textbf{3. Doppelstunde} (Std. 5-6 von 12) der Sequenz ``\reihenthema''.

\subsection{Kompetenzziele (Rahmenplan MV)}
\begin{itemize}
    \item \textbf{Sprachlich:} Possessivbegleiter (mi, tu, su) in Singular und Plural verwenden
    \item \textbf{Kommunikativ:} Über die eigene Familie sprechen und nach der Familie anderer fragen
    \item \textbf{Interkulturell:} Familien in Spanien (Bedeutung der familia extendida)
\end{itemize}

\subsection{Legitimation}
\textbf{Rahmenplan Spanisch MV:}
\begin{itemize}
    \item Verfügen über sprachliche Mittel: Grammatik -- Possessivbegleiter
    \item Sprechen: Über Personen und Beziehungen sprechen
    \item Themenbereich: yo y mi entorno -- Familie
\end{itemize}

% ============================================================================
% 4. STUNDENZIELE
% ============================================================================
\section{Stundenziele}

\subsection{Grobziel}
Die SuS erweitern ihre \textbf{grammatische Kompetenz}, indem sie die Possessivbegleiter (mi/mis, tu/tus, su/sus) verstehen und anwenden, um \textbf{über ihre Familie zu sprechen}.

\subsection{Feinziele (operationalisiert)}
\begin{tabular}{|c|p{8cm}|c|c|}
\hline
\textbf{Nr.} & \textbf{Die SuS können...} & \textbf{AFB} & \textbf{Diff.} \\
\hline
FZ1 & die Possessivbegleiter mi, tu, su benennen und ins Deutsche übersetzen & I & alle \\
\hline
FZ2 & die Pluralformen mis, tus, sus korrekt bilden und verwenden & I/II & alle \\
\hline
FZ3 & Familienmitglieder mit Possessivbegleitern beschreiben (Mi madre se llama...) & II & alle \\
\hline
FZ4 & Fragen über die Familie stellen und beantworten (¿Cómo se llama tu...?) & II & \Std{}/\E{} \\
\hline
FZ5 & die Regel zur Kongruenz (Bezugswort statt Besitzer) erklären und anwenden & II/III & \E{} \\
\hline
\end{tabular}

% ============================================================================
% 5. METHODISCHE ANALYSE
% ============================================================================
\section{Methodische Analyse}

\subsection{Medien}
\begin{itemize}
    \item Tafel/Whiteboard
    \item Lehrwerk: \lehrwerk, \lehrwerkseite
    \item Arbeitsblatt: AB-02c.pdf
    \item Bildkarten: Familienstammbaum
    \item Optional: LP-02c.html (digital)
    \item Familien-Flashcards (mi madre, tu padre, su hermano...)
\end{itemize}

\subsection{Sozialformen}
\begin{itemize}
    \item Plenum: Einführung Possessivbegleiter, Regelbildung
    \item Partnerarbeit: Dialog-Übung (Fragen nach Familie)
    \item Einzelarbeit: AB-Aufgaben
\end{itemize}

% ============================================================================
% 6. VERLAUFSPLANUNG
% ============================================================================
\section{Verlaufsplanung}

\textbf{1. Stunde (45 Min.) -- Schwerpunkt: mi/tu/su (Singular)}

\begin{longtable}{|p{1cm}|p{1.5cm}|p{4cm}|p{3cm}|p{2cm}|p{2cm}|}
\hline
\textbf{Zeit} & \textbf{Phase} & \textbf{Lehrerhandeln} & \textbf{Schülerhandeln} & \textbf{SF} & \textbf{Medien} \\
\hline
\endhead

0--5 & Einstieg & Familienfoto zeigen: ``Mi familia'' -- ``¿Y tu familia?'' & Vermutungen äußern & Plenum & Bild \\
\hline

5--15 & Input & Possessiv einführen: mi madre, tu padre, su hermano; Tafelanschrieb & Mitschreiben, Nachsprechen & Plenum & Tafel \\
\hline

15--25 & Übung 1 & Bildkarten: ``¿mi, tu oder su?'' -- L zeigt Familienmitglieder & Chorisch antworten: ``mi madre'', ``tu hermano'' & Plenum & Karten \\
\hline

25--35 & Übung 2 & AB-02c Aufg. 1 erklären & Lückentext ausfüllen & EA & AB \\
\hline

35--40 & Sicherung & Lösungen besprechen, Fragen klären & Vergleichen, korrigieren & Plenum & Tafel \\
\hline

40--45 & Überleitung & Plural einführen: mi $\rightarrow$ mis, tu $\rightarrow$ tus, su $\rightarrow$ sus & Notieren, Beispiele nennen & Plenum & Tafel \\
\hline
\end{longtable}

\vspace{1em}

\textbf{2. Stunde (45 Min.) -- Schwerpunkt: Plural + Anwendung}

\begin{longtable}{|p{1cm}|p{1.5cm}|p{4cm}|p{3cm}|p{2cm}|p{2cm}|}
\hline
\textbf{Zeit} & \textbf{Phase} & \textbf{Lehrerhandeln} & \textbf{Schülerhandeln} & \textbf{SF} & \textbf{Medien} \\
\hline
\endhead

0--5 & Aktivierung & Quiz: ``mi oder mis?'' -- hermano vs. hermanos & Schnell antworten & Plenum & -- \\
\hline

5--15 & Übung 3 & AB-02c Aufg. 2: Singular $\rightarrow$ Plural umformen & Sätze transformieren & EA & AB \\
\hline

15--25 & Anwendung & AB-02c Aufg. 3: Familie beschreiben ``Mi madre se llama...'' & Sätze schreiben & EA & AB \\
\hline

25--35 & Dialog & PA: ``¿Cómo se llama tu hermano?'' -- ``Mi hermano se llama...'' & Dialoge führen & PA & AB \\
\hline

35--40 & Sicherung & Merkkasten AB-02c gemeinsam ausfüllen & Übertragen, vergleichen & EA/PL & AB \\
\hline

40--45 & Abschluss & Zusammenfassung, HA erklären & Aufschreiben & Plenum & -- \\
\hline
\end{longtable}

\textbf{Legende:} EA = Einzelarbeit, PA = Partnerarbeit, PL = Plenum, SF = Sozialform

% ============================================================================
% 7. ERWARTETE SCHWIERIGKEITEN
% ============================================================================
\section{Erwartete Schwierigkeiten}

\begin{tabular}{|p{5cm}|p{8cm}|}
\hline
\textbf{Schwierigkeit} & \textbf{Reaktion} \\
\hline
Verwechslung mi/mis & Regel betonen: mi + Singular, mis + Plural; Visualisierung: mi libro vs. mis libros \\
\hline
Genusfehler (*mia madre) & Im Spanischen kein Unterschied: mi madre = mi padre (nur Numerus!) \\
\hline
su-Mehrdeutigkeit & Kontext klären: ``su madre'' = seine/ihre Mutter; im Dialog nachfragen \\
\hline
Transfer aus Deutsch & Bewusstmachung: mein/meine (Genus) vs. mi (kein Genus), nur mis (Plural) \\
\hline
Aussprache tu vs. tú & Betonung: tu (Possessiv, unbetont) vs. tú (Pronomen, betont) \\
\hline
\end{tabular}

% ============================================================================
% 8. HAUSAUFGABE
% ============================================================================
\section{Hausaufgabe}

\begin{itemize}
    \item Lehrwerk S. 65, Übung 5 (Possessivbegleiter einsetzen)
    \item Vokabeln lernen: Possessivbegleiter + 5 Familienwörter
    \item Optional: LP-02c.html bearbeiten (interaktive Übungen)
    \item Kreativ: Stammbaum der eigenen Familie zeichnen und beschriften
\end{itemize}

% ============================================================================
% 9. LÖSUNGEN
% ============================================================================
\section{Lösungen AB-02c}

\textbf{Merkkasten: Possessivbegleiter}

\begin{tabular}{lll}
yo: & mi (Sg.) & mis (Pl.) \\
tú: & tu (Sg.) & tus (Pl.) \\
él/ella/usted: & su (Sg.) & sus (Pl.) \\
\end{tabular}

\vspace{1em}

\textbf{Aufgabe 1: mi/tu/su einsetzen} (4 Punkte)
\begin{enumerate}[label=\alph*)]
    \item \textbf{Mi} madre se llama Carmen.
    \item ¿Cómo se llama \textbf{tu} padre?
    \item \textbf{Su} hermano tiene 15 años.
    \item ¿Cuántos años tiene \textbf{tu} abuela?
\end{enumerate}

\vspace{1em}

\textbf{Aufgabe 2: Singular $\rightarrow$ Plural} (4 Punkte)
\begin{enumerate}[label=\alph*)]
    \item mi hermano $\rightarrow$ \textbf{mis hermanos}
    \item tu abuelo $\rightarrow$ \textbf{tus abuelos}
    \item su prima $\rightarrow$ \textbf{sus primas}
    \item mi tío $\rightarrow$ \textbf{mis tíos}
\end{enumerate}

\vspace{1em}

\textbf{Aufgabe 3: Familie beschreiben} (6 Punkte)

Beispielantworten:
\begin{itemize}
    \item Mi madre se llama [Name]. Tiene [Alter] años.
    \item Mi padre se llama [Name].
    \item Mi hermano/hermana se llama [Name].
    \item Mis abuelos se llaman [Name] y [Name].
\end{itemize}

\textit{Bewertung: Je 2P pro korrektem Satz (Possessiv + llamarse + Info)}

\vspace{1em}

\textbf{Aufgabe 4: Dialog} (6 Punkte)

Beispieldialog:
\begin{itemize}
    \item A: ¿Cómo se llama \textbf{tu} hermano?
    \item B: \textbf{Mi} hermano se llama Pablo.
    \item A: ¿Cuántos años tiene?
    \item B: Tiene doce años.
    \item A: ¿Y \textbf{tu} madre?
    \item B: \textbf{Mi} madre se llama Ana.
\end{itemize}

\textit{Bewertung: Fragen korrekt (2P), Antworten mit Possessiv (2P), Grammatik (2P)}

% ============================================================================
% 10. DIDAKTIK-SCORE
% ============================================================================
\newpage
\section{Didaktik-Score}

\begin{scorebox}
\textbf{Bewertung der didaktischen Qualität nach 10 Dimensionen}\\
\textbf{Ziel-Score: $\geq$ 9.0 (Premium-Qualität)}

\vspace{1em}
\begin{tabular}{|c|l|c|c|p{4.5cm}|}
\hline
\textbf{\#} & \textbf{Dimension} & \textbf{Gew.} & \textbf{Score} & \textbf{Notiz} \\
\hline
1 & Differenzierung & 15\% & 9/10 & [B] Lückentext, [S] Sätze bilden, [E] freie Beschreibung \\
\hline
2 & Taxonomie (AFB) & 10\% & 9/10 & AFB I: 35\%, AFB II: 50\%, AFB III: 15\% \\
\hline
3 & Lernziel-Alignment & 10\% & 10/10 & RP MV: Possessivbegleiter, Familie \\
\hline
4 & Aktivierung & 10\% & 9/10 & Chorisch, PA-Dialog, Quiz \\
\hline
5 & Scaffolding & 10\% & 9/10 & Singular $\rightarrow$ Plural $\rightarrow$ freie Anwendung \\
\hline
6 & Feedback-Qualität & 10\% & 9/10 & Elaboriertes Feedback in LP + Musterlösung \\
\hline
7 & Sprachliche Klarheit & 10\% & 10/10 & Altersgerecht, kurze Sätze, Tabellen \\
\hline
8 & Motivation \& Relevanz & 10\% & 9/10 & Familie = Alltagsrelevanz, persönlicher Bezug \\
\hline
9 & Inklusion (UDL) & 10\% & 9/10 & Visuell (Tabelle) + auditiv + kinästhetisch \\
\hline
10 & Praktikabilität & 5\% & 9/10 & 2x45 Min realistisch, Material vorhanden \\
\hline
\multicolumn{3}{|r|}{\textbf{GESAMT:}} & \textbf{9.2/10} & \textcolor{successgreen}{\textbf{FREIGABE}} \\
\hline
\end{tabular}
\end{scorebox}

\end{document}
